\documentclass[11pt,a4paper]{article}
\usepackage[margin=1.8cm]{geometry}
\usepackage{amsmath,amssymb}
\usepackage{booktabs,array,tabularx}
\usepackage{tcolorbox}
\usepackage{microtype}
\usepackage{parskip}
\usepackage{enumitem}
\usepackage{xcolor}   % used only for gray shading — no hue colours
\tcbuselibrary{skins,breakable}

% ── All boxes: black/white/gray only ─────────────────────────────────────────
\tcbset{
  defn/.style={colback=gray!8,colframe=black,fonttitle=\bfseries\small,
    title=#1,left=4pt,right=4pt,top=3pt,bottom=3pt,boxrule=0.7pt,arc=2pt},
  bold/.style={colback=white,colframe=black,fonttitle=\bfseries\small,
    title=#1,left=4pt,right=4pt,top=3pt,bottom=3pt,boxrule=1.2pt,arc=2pt},
  exam/.style={colback=gray!5,colframe=black,fonttitle=\bfseries\small,
    title=#1,left=4pt,right=4pt,top=3pt,bottom=3pt,
    boxrule=0.5pt,arc=2pt,borderline west={2pt}{0pt}{black}},
  note/.style={colback=white,colframe=gray!60,fonttitle=\bfseries\small,
    title=#1,left=4pt,right=4pt,top=3pt,bottom=3pt,boxrule=0.5pt,arc=2pt}
}

\setlist[itemize]{noitemsep,topsep=3pt,parsep=0pt}
\setlist[enumerate]{noitemsep,topsep=3pt,parsep=0pt}

\newcommand{\E}{\mathbb{E}}
\newcommand{\Qt}{Q_t}
\newcommand{\qstar}{q^*(a)}

\begin{document}

% ── Title ────────────────────────────────────────────────────────────────────
\begin{center}
  {\large\bfseries DRL Sessions 2 \& 3 — Multi-Armed Bandits (Exam Handout)}\\[2pt]
  {\small Expected Learnings: MAB formulation · Action-value methods · $\varepsilon$-greedy ·
  Incremental updates · Non-stationarity · Optimistic IV · UCB · Bandit Gradient · Contextual Bandits}
\end{center}
\hrule\vspace{6pt}

% ════════════════════════════════════════════════════════════════════════════
\section*{1 · The \(k\)-Armed Bandit Problem}
% ════════════════════════════════════════════════════════════════════════════

\begin{tcolorbox}[defn={Formal Definition}]
At each time step $t$ you choose one of $k$ actions (arms). You receive a numerical reward
drawn from the action's \textbf{stationary} probability distribution.
\textbf{Objective:} maximise \emph{expected total reward} over some time horizon.
\end{tcolorbox}

\textbf{Three defining features}
\begin{enumerate}
  \item \textbf{Unknown rewards} — the true value $q^*(a) = \E[R_t \mid A_t = a]$ is not known upfront.
  \item \textbf{Repeated choice} — you face the same problem over and over and can learn from experience.
  \item \textbf{Observable outcome} — after choosing $a$ you observe the reward $R_t$.
\end{enumerate}

\textbf{Why we study MAB:} It is the simplest setting where the \emph{exploration–exploitation trade-off}
appears. Every concept (value estimates, greedy selection, $\varepsilon$-greedy, UCB) reappears
in full MDPs.

% ════════════════════════════════════════════════════════════════════════════
\section*{2 · Action-Value Methods (Sample Average)}
% ════════════════════════════════════════════════════════════════════════════

\begin{tcolorbox}[defn={Estimated value of action $a$ at time $t$}]
\[
  Q_t(a) \;=\; \frac{\text{sum of rewards when } a \text{ was selected before } t}
                    {\text{number of times } a \text{ was selected before } t}
           \;=\; \frac{\sum_{i=1}^{t-1} R_i \cdot \mathbf{1}_{A_i=a}}
                      {\sum_{i=1}^{t-1} \mathbf{1}_{A_i=a}}
\]
By the Law of Large Numbers, $Q_t(a) \to q^*(a)$ as the count $\to \infty$.
\end{tcolorbox}

\subsection*{Greedy Action Selection}
\[
  A_t^* \;=\; \arg\max_a Q_t(a)
\]
\textit{Always} exploiting — never discovers potentially better actions. Optimal only if initial estimates
are already correct.

% ════════════════════════════════════════════════════════════════════════════
\section*{3 · Exploration vs.\ Exploitation — \(\varepsilon\)-Greedy}
% ════════════════════════════════════════════════════════════════════════════

\begin{tcolorbox}[bold={$\varepsilon$-Greedy Policy (Algorithm)}]
\begin{verbatim}
epsilon = 0.1            # small exploration rate
def get_action():
    if random() > epsilon:        # prob (1 - epsilon)
        return argmax Q(a)        # EXPLOIT: greedy
    else:                         # prob epsilon
        return random.choice(A)   # EXPLORE: uniform random
\end{verbatim}
\end{tcolorbox}

\textbf{Probabilities for a specific action $a$ (greedy action):}
\begin{align*}
  P(\text{greedy action selected}) &= \underbrace{(1-\varepsilon)}_{\text{exploit branch}}
    + \underbrace{\frac{\varepsilon}{k}}_{\text{explore branch picks greedy by chance}} \\[4pt]
  P(\text{non-greedy action } a_i \text{ selected}) &= \frac{\varepsilon}{k}
\end{align*}

\begin{tcolorbox}[exam={Worked Example 1 — from slides (must know)}]
\textbf{Q:} In $\varepsilon$-greedy with $k = 2$ actions and $\varepsilon = 0.5$,
what is the probability that the greedy action is selected?

\medskip
\textbf{Method (law of total probability):}
\begin{align*}
  P(\text{greedy}) &= P(\text{greedy} \mid \text{greedy branch})\,P(\text{greedy branch})\\
                   &\quad + P(\text{greedy} \mid \text{explore branch})\,P(\text{explore branch})\\
                   &= 1 \times (1 - 0.5) + \tfrac{1}{2} \times 0.5\\
                   &= 0.5 + 0.25 \;=\; \mathbf{0.75}
\end{align*}

\textbf{General formula check:} $(1 - \varepsilon) + \varepsilon/k = 0.5 + 0.5/2 = 0.75\checkmark$
\end{tcolorbox}

\subsection*{Effect of $\varepsilon$ on performance}
\begin{center}
\begin{tabular}{p{2.8cm}p{5.5cm}p{5.5cm}}
\toprule
$\varepsilon$ & Behaviour & Best when \\
\midrule
$0$ (pure greedy) & Never explores; gets stuck at a sub-optimal arm & Rewards deterministic; initial estimates already good \\
Small (0.01--0.1) & Mostly exploits; finds optimum slowly but surely & Stationary, noisy rewards \\
Large (0.5) & Explores heavily; slow to exploit best arm & High variance rewards or non-stationary setting \\
\bottomrule
\end{tabular}
\end{center}

\textbf{10-armed testbed conclusion:} $\varepsilon = 0.1$ finds the best arm faster early on;
$\varepsilon = 0.01$ eventually performs better as its estimates become more accurate.

% ════════════════════════════════════════════════════════════════════════════
\section*{4 · Incremental Implementation (Session 3)}
% ════════════════════════════════════════════════════════════════════════════

\begin{tcolorbox}[defn={Incremental update rule}]
Given $Q_n$ (estimate after $n-1$ rewards) and the $n$-th reward $R_n$:
\[
  Q_{n+1} \;=\; Q_n + \frac{1}{n}\bigl[R_n - Q_n\bigr]
\]
\textbf{General form (step size $\alpha$):}
\[
  Q_{n+1} \;=\; Q_n + \alpha\bigl[R_n - Q_n\bigr]
\]
$[R_n - Q_n]$ is the \textbf{error} (target $-$ estimate). Multiply by step size; add to old estimate.
No need to store all past rewards.
\end{tcolorbox}

\textbf{Constant vs.\ Varying Step Size}
\begin{itemize}
  \item $\alpha = 1/n$ (sample average): gives equal weight to all past rewards. Good for \textbf{stationary} problems.
  \item $\alpha = $ constant: gives more weight to \emph{recent} rewards ($\alpha(1-\alpha)^{n-k}$ weight on reward $k$ steps ago).
    Good for \textbf{non-stationary} problems.
\end{itemize}

% ════════════════════════════════════════════════════════════════════════════
\section*{5 · Non-Stationarity}
% ════════════════════════════════════════════════════════════════════════════

\textbf{Problem:} most real RL environments change over time.
$\varepsilon$-greedy with sample average ($\alpha = 1/n$) converges to the true value of the
\emph{original} distribution — not the current one.

\begin{tcolorbox}[defn={Exponential recency-weighted average ($\alpha$ constant)}]
\[
  Q_{n+1} = (1-\alpha)^n Q_1 + \sum_{i=1}^{n} \alpha(1-\alpha)^{n-i} R_i
\]
Recent rewards get exponentially more weight. This is the standard approach for non-stationary bandits.
\end{tcolorbox}

% ════════════════════════════════════════════════════════════════════════════
\section*{6 · Optimistic Initial Values}
% ════════════════════════════════════════════════════════════════════════════

\textbf{Idea:} Set $Q_1(a) = +5$ (much higher than any real reward, e.g.\ max real reward $\approx 1$).
Every action is initially over-estimated.

\textbf{Effect:} Whichever action is tried first yields a real reward $< 5$.
The agent is ``disappointed'' and tries other actions. All actions are explored early.

\textbf{Caution (exam point):}
\begin{itemize}
  \item Works well for \textbf{stationary} problems only.
  \item In \textbf{non-stationary} settings it fails: once the initial exploration burst is over,
    the agent stops exploring, even if action values have since changed.
\end{itemize}

% ════════════════════════════════════════════════════════════════════════════
\section*{7 · Upper Confidence Bound (UCB)}
% ════════════════════════════════════════════════════════════════════════════

\begin{tcolorbox}[defn={UCB Action Selection}]
\[
  A_t = \arg\max_a \left[ Q_t(a) + c\sqrt{\frac{\ln t}{N_t(a)}} \right]
\]
\textbf{Each term:}
\begin{itemize}
  \item $Q_t(a)$ — current value estimate of action $a$
  \item $N_t(a)$ — how many times $a$ has been selected so far
  \item $c\sqrt{\ln t / N_t(a)}$ — \textbf{uncertainty bonus}: large when $a$ is rarely tried; shrinks as $a$ is tried more
  \item $c$ — confidence level; controls how optimistic we are about uncertain actions
\end{itemize}
\end{tcolorbox}

\textbf{Intuition:} UCB explores \emph{intelligently} — it prefers actions that are
\emph{either high-valued or rarely tried}, not random ones.
As $t$ grows, $\ln t / N_t(a) \to 0$ for frequently tried arms — they are selected less often over time.

\textbf{Limitation:} UCB is harder to generalise beyond the bandit setting to full MDPs (unlike $\varepsilon$-greedy).

% ════════════════════════════════════════════════════════════════════════════
\section*{8 · Bandit Gradient Algorithm (High Level)}
% ════════════════════════════════════════════════════════════════════════════

Instead of estimating action \emph{values}, learn action \emph{preferences} $H_t(a)$ directly.

\textbf{Action selection — Softmax:}
\[
  \pi_t(a) \;=\; \frac{e^{H_t(a)}}{\sum_{b=1}^k e^{H_t(b)}}
\]

\textbf{Preference update after selecting $A_t$ and receiving $R_t$:}
\[
  H_{t+1}(A_t) = H_t(A_t) + \alpha(R_t - \bar{R}_t)(1 - \pi_t(A_t))
\]
\[
  H_{t+1}(a) = H_t(a) - \alpha(R_t - \bar{R}_t)\pi_t(a) \quad \forall\, a \neq A_t
\]
where $\bar{R}_t$ is the running average reward (the baseline).

\textbf{Key idea:} If $R_t > \bar{R}_t$ (better than average), increase preference for $A_t$;
decrease for all others. Proofs excluded from exam scope.

% ════════════════════════════════════════════════════════════════════════════
\section*{9 · Contextual Bandits (Associative Tasks)}
% ════════════════════════════════════════════════════════════════════════════

\begin{tcolorbox}[note={What are Contextual Bandits?}]
In the standard bandit, the same $k$ arms exist in every step.
In a \textbf{contextual bandit}, each step has a \textbf{context} (situation/state) $s$,
and the best action can differ by context.

\textbf{Example:} A news recommender. Context = user profile + time of day.
Best article to recommend changes with the context.

\textbf{Policy:} a mapping from context $\to$ best action (unlike standard MAB which maps nothing $\to$ action).

\textbf{Why special approaches?} If context switches are rapid, optimistic initial values and simple
averaging fail. Need methods that generalise across contexts (e.g.\ linear regression on features,
neural network policies).

\textbf{Relation to full RL:} Contextual bandits are the bridge from MAB to MDP.
In an MDP, your action changes the \emph{next} context — that dependency is absent in contextual bandits.
\end{tcolorbox}

% ════════════════════════════════════════════════════════════════════════════
\section*{10 · Exam Worked Examples}
% ════════════════════════════════════════════════════════════════════════════

\begin{tcolorbox}[exam={Ex-2 — Which steps definitely/possibly used the $\varepsilon$ case?}]
\textbf{Setup:} $k=4$, $\varepsilon$-greedy, sample-average, $Q_1(a)=0$ for all $a$.\\
\textbf{Sequence:} $A_1=1,R_1=1;\; A_2=2,R_2=1;\; A_3=2,R_3=2;\; A_4=2,R_4=2;\; A_5=3,R_5=0$

\medskip
Compute $Q$ before each step, then check if the action was greedy:
\begin{center}
\small
\begin{tabular}{ccccccc}
\toprule
Step & $Q(1)$ & $Q(2)$ & $Q(3)$ & $Q(4)$ & Greedy? & $\varepsilon$ case? \\
\midrule
$t=1$ & 0 & 0 & 0 & 0 & Tie (any) & \textbf{Possibly} (tie $\to$ any arm is ``greedy'') \\
$t=2$ & 1.0 & 0 & 0 & 0 & $A^*=1$ & \textbf{Definitely} (chose 2, not greedy arm 1) \\
$t=3$ & 1.0 & 1.0 & 0 & 0 & Tie (1 or 2) & \textbf{Possibly} (arm 2 is tied for greedy) \\
$t=4$ & 1.0 & 1.5 & 0 & 0 & $A^*=2$ & Possibly (arm 2 is greedy; choosing it is fine) \\
$t=5$ & 1.0 & 1.67 & 0 & 0 & $A^*=2$ & \textbf{Definitely} (chose 3, not greedy arm 2) \\
\bottomrule
\end{tabular}
\end{center}
\textbf{Answer:} $\varepsilon$ case \emph{definitely} occurred at $t=2$ and $t=5$.
It \emph{possibly} occurred at $t=1$ (tie), $t=3$ (tie), $t=4$ (arm 2 was greedy so no $\varepsilon$ needed).
\end{tcolorbox}

\begin{tcolorbox}[exam={Summary — All 4 Methods at a Glance}]
\begin{tabular}{p{3.2cm}p{5cm}p{5.5cm}}
\toprule
\textbf{Method} & \textbf{Key idea} & \textbf{Best for} \\
\midrule
Greedy & Always pick highest $Q_t(a)$ & Only if estimates are already good \\
$\varepsilon$-Greedy & Exploit $(1-\varepsilon)$, random explore $\varepsilon$ & General-purpose; easy to tune \\
Optimistic IV & Start estimates high; forces exploration early & Stationary problems only \\
UCB & Exploit + uncertainty bonus & Stationary; smarter exploration than random \\
Gradient/Softmax & Learn preferences, not values & Policy-based; when value estimation is hard \\
\bottomrule
\end{tabular}
\end{tcolorbox}

\vspace{6pt}
\hrule
\vspace{3pt}
{\footnotesize DRL EC3 Exam Handout — Sessions 2 \& 3 (Multi-Armed Bandits). Ref: Sutton \& Barto Ch.~2.}

\end{document}
